\documentclass[10pt]{article}
\usepackage[T1]{fontenc}
\usepackage{lmodern}
\usepackage[margin=22mm]{geometry}
\usepackage{xcolor}
\usepackage{microtype}
\usepackage{enumitem}
\usepackage{hyperref}
\usepackage{xurl}
\usepackage{forest}
\usetikzlibrary{arrows.meta}
\usepackage{pdflscape}
\usepackage[most]{tcolorbox}
\usepackage{titlesec}
\usepackage{fancyhdr}
\definecolor{Navy}{HTML}{17324D}
\definecolor{Teal}{HTML}{176B6B}
\definecolor{PaleBlue}{HTML}{EAF2F8}
\definecolor{FileFill}{HTML}{F7F9FA}
\definecolor{DescText}{HTML}{334E5C}
\definecolor{TreeLine}{HTML}{607D8B}
\definecolor{Warn}{HTML}{8A4B08}
\hypersetup{colorlinks=true,linkcolor=Teal,urlcolor=Teal,pdfborder={0 0 0}}
\titleformat{\section}{\Large\bfseries\color{Navy}}{}{0pt}{}[\color{Teal}\titlerule]
\titleformat{\subsection}{\large\bfseries\color{Teal}}{}{0pt}{}
\setlist[itemize]{leftmargin=*,itemsep=.45em,topsep=.45em}
\setlist[enumerate]{leftmargin=*,itemsep=.45em,topsep=.45em}
\setlength{\parindent}{0pt}
\setlength{\parskip}{.65em}
\newcommand{\code}[1]{\texttt{#1}}
\forestset{
 rootnode/.style={draw=Navy,fill=Navy,text=white,rounded corners=2pt,inner sep=5pt},
 filenode/.style={draw=TreeLine,fill=FileFill,rounded corners=2pt,inner sep=3.2pt},
}
\pagestyle{fancy}\fancyhf{}\fancyhead[L]{\small\color{DescText}Repository Structure Inventory}\fancyhead[R]{\small\color{DescText}\leftmark}\fancyfoot[C]{\thepage}
\setlength{\headheight}{14pt}
\newtcolorbox{verificationbox}{breakable,colback=orange!5,colframe=Warn,title={Verification status},fonttitle=\bfseries}
\begin{document}
\hypersetup{pageanchor=false}
\begin{titlepage}
\pagecolor{PaleBlue}
\color{Navy}
\vspace*{0.17\textheight}
{\Huge\bfseries Grade 11 — Term 2 Structure\par}
\vspace{1em}
{\Large Repository resource inventory\par}
\vfill
{\large Authoritative conversion of \code{STRUCTURE.md}\par}
\end{titlepage}
\nopagecolor
\hypersetup{pageanchor=true}
\pagenumbering{arabic}
\tableofcontents
\clearpage

\section{Overview}

This directory contains the LaTeX sources for a Grade 11 statistics unit on confidence-interval planning and inference in economic and financial contexts. Its four collections separate worked practice, formal learning evidence and assessments, criterion-specific catch-up work, and supplementary lesson material. Student-facing question sheets generally have a corresponding teacher-facing worked-solution source; no compiled PDFs, datasets, web assets, or build configuration are stored here.

\code{STRUCTURE.md} is this inventory document and is intentionally omitted from the inventory counts and tree below.

\section{Directory tree}

\begin{tcolorbox}[colback=PaleBlue,colframe=Teal,title={Reading the diagrams}]
The directory tree is divided into one landscape diagram per top-level directory so that every documented entry and description remains legible. Directory nodes use a dark fill; file nodes use a light fill, with descriptions shown in a secondary text color. The inventory root is \code{.}.
\end{tcolorbox}
\begin{landscape}
\subsection{\code{1-practice\_activities/}}
\begin{center}
\begin{forest}
for tree={grow=east, parent anchor=east, child anchor=west, anchor=west, edge={-Latex,draw=TreeLine}, l sep=6mm, s sep=1.2mm, font=\sffamily\fontsize{6.2}{7.2}\selectfont},
[{\textbf{1-practice\_activities/}\\\textcolor{DescText}{Worked lesson practice organized by lesson and assessment criterion.}}, rootnode, text width=6.0cm
[{\code{G11\_T2\_L2\_C6\_practice.tex}\\\textcolor{DescText}{Worked C6 practice on required and additional sample sizes for means with known population spread.}}, filenode, text width=16.0cm]
[{\code{G11\_T2\_L2\_C7\_practice.tex}\\\textcolor{DescText}{Worked C7 practice constructing and interpreting confidence intervals for proportions.}}, filenode, text width=16.0cm]
[{\code{G11\_T2\_L3\_C10\_practice.tex}\\\textcolor{DescText}{Worked C10 hypothesis-test conclusions for mean parameters, including Type I and II errors.}}, filenode, text width=16.0cm]
[{\code{G11\_T2\_L3\_C8\_practice.tex}\\\textcolor{DescText}{Worked C8 small-sample t-interval examples for means with unknown variance.}}, filenode, text width=16.0cm]
[{\code{G11\_T2\_L3\_C9\_practice.tex}\\\textcolor{DescText}{Worked C9 comparisons of t- and z-intervals for means when variance is unknown.}}, filenode, text width=16.0cm]
]
\end{forest}
\end{center}
\end{landscape}
\begin{landscape}
\subsection{\code{2-learning\_evidences/}}
\begin{center}
\begin{forest}
for tree={grow=east, parent anchor=east, child anchor=west, anchor=west, edge={-Latex,draw=TreeLine}, l sep=6mm, s sep=1.2mm, font=\sffamily\fontsize{6.2}{7.2}\selectfont},
[{\textbf{2-learning\_evidences/}\\\textcolor{DescText}{Student assessments, project briefs, practice assessments, and teacher solutions.}}, rootnode, text width=6.0cm
[{\code{G11\_T2\_L1\_exam.tex}\\\textcolor{DescText}{Student-facing Lesson 1 exam on sample summaries and mean intervals with known population standard deviation.}}, filenode, text width=16.0cm]
[{\code{G11\_T2\_L1\_exam\_practice.tex}\\\textcolor{DescText}{Student-facing two-problem practice exam using grouped loading times and raw cafeteria-demand data.}}, filenode, text width=16.0cm]
[{\code{G11\_T2\_L1\_exam\_practice\_solution.tex}\\\textcolor{DescText}{Teacher-facing worked solutions to the Lesson 1 practice exam.}}, filenode, text width=16.0cm]
[{\code{G11\_T2\_L1\_exam\_solution.tex}\\\textcolor{DescText}{Teacher-facing worked solutions to the Lesson 1 exam's subsidy, order-rate, and loading-time problems.}}, filenode, text width=16.0cm]
[{\code{G11\_T2\_L2\_C6\_activity.tex}\\\textcolor{DescText}{Student-facing C6 activity with six mean sample-size and margin-of-error planning problems.}}, filenode, text width=16.0cm]
[{\code{G11\_T2\_L2\_C6\_activity\_solution.tex}\\\textcolor{DescText}{Teacher-facing worked solutions to all six C6 activity problems.}}, filenode, text width=16.0cm]
[{\code{G11\_T2\_L2\_C7\_activity.tex}\\\textcolor{DescText}{Student-facing C7 activity assessing proportion intervals in five applied contexts plus follow-up problems.}}, filenode, text width=16.0cm]
[{\code{G11\_T2\_L2\_C7\_activity\_solution.tex}\\\textcolor{DescText}{Teacher-facing worked solutions for the C7 activity's contextual and numbered problems.}}, filenode, text width=16.0cm]
[{\code{G11\_T2\_L2\_midterm.tex}\\\textcolor{DescText}{Student-facing midterm covering sample statistics, parameters, confidence intervals, interpretation, and sample size.}}, filenode, text width=16.0cm]
[{\code{G11\_T2\_L2\_midterm\_solution.tex}\\\textcolor{DescText}{Teacher-facing fully worked solutions to the standard midterm.}}, filenode, text width=16.0cm]
[{\code{G11\_T2\_L2\_mock\_midterm.tex}\\\textcolor{DescText}{Student-facing standard mock midterm on mean/proportion intervals and sample-size planning.}}, filenode, text width=16.0cm]
[{\code{G11\_T2\_L2\_mock\_midterm\_hard.tex}\\\textcolor{DescText}{Student-facing harder mock with more demanding multi-step inference and planning tasks.}}, filenode, text width=16.0cm]
[{\code{G11\_T2\_L2\_mock\_midterm\_hard\_solution.tex}\\\textcolor{DescText}{Teacher-facing worked solutions to the harder mock midterm.}}, filenode, text width=16.0cm]
[{\code{G11\_T2\_L2\_mock\_midterm\_solution.tex}\\\textcolor{DescText}{Teacher-facing worked solutions to the standard mock midterm.}}, filenode, text width=16.0cm]
[{\code{G11\_T2\_L2\_reference\_midterm.tex}\\\textcolor{DescText}{Worked reference midterm for confidence and interval estimation; it is labeled as solutions, not a blank student paper.}}, filenode, text width=16.0cm]
[{\code{G11\_T2\_L3\_proyect.tex}\\\textcolor{DescText}{Student-facing project brief for presenting and defending an AI-generated confidence-interval web page.}}, filenode, text width=16.0cm]
[{\code{G11\_T2\_L3\_proyect\_Excel.tex}\\\textcolor{DescText}{Alternative student project brief for an Excel confidence-interval tool with CSV input, dashboards, and simple macros.}}, filenode, text width=16.0cm]
[{\code{G11\_T2\_L3\_proyect\_problems.tex}\\\textcolor{DescText}{Student-facing set of ten t-interval calculation, interpretation, and decision problems for the project unit.}}, filenode, text width=16.0cm]
[{\code{G11\_T2\_L3\_proyect\_problems\_solutions.tex}\\\textcolor{DescText}{Teacher-facing detailed and concise solutions to the ten project-unit problems.}}, filenode, text width=16.0cm]
]
\end{forest}
\end{center}
\end{landscape}
\begin{landscape}
\subsection{\code{3-catch\_ups/}}
\begin{center}
\begin{forest}
for tree={grow=east, parent anchor=east, child anchor=west, anchor=west, edge={-Latex,draw=TreeLine}, l sep=6mm, s sep=1.2mm, font=\sffamily\fontsize{6.2}{7.2}\selectfont},
[{\textbf{3-catch\_ups/}\\\textcolor{DescText}{Criterion-by-criterion remediation sheets and their teacher solution companions.}}, rootnode, text width=6.0cm
[{\code{G11\_T2\_C1\_catchUp.tex}\\\textcolor{DescText}{Student-facing C1 practice computing sample means and standard deviations in ten finance contexts.}}, filenode, text width=16.0cm]
[{\code{G11\_T2\_C1\_catchUp\_solution.tex}\\\textcolor{DescText}{Teacher-facing worked solutions to the ten C1 problems.}}, filenode, text width=16.0cm]
[{\code{G11\_T2\_C2\_catchUp.tex}\\\textcolor{DescText}{Student-facing C2 practice distinguishing sample statistics from population parameters.}}, filenode, text width=16.0cm]
[{\code{G11\_T2\_C2\_catchUp\_solution.tex}\\\textcolor{DescText}{Teacher-facing detailed and minimal solutions to the ten C2 classification problems.}}, filenode, text width=16.0cm]
[{\code{G11\_T2\_C3\_catchUp.tex}\\\textcolor{DescText}{Student-facing C3 practice comparing point and interval estimates and their decision risk.}}, filenode, text width=16.0cm]
[{\code{G11\_T2\_C3\_catchUp\_solution.tex}\\\textcolor{DescText}{Teacher-facing contextual explanations for the ten C3 problems.}}, filenode, text width=16.0cm]
[{\code{G11\_T2\_C4\_catchUp.tex}\\\textcolor{DescText}{Student-facing C4 practice interpreting interval bounds, confidence, plausibility, and comparisons.}}, filenode, text width=16.0cm]
[{\code{G11\_T2\_C4\_catchUp\_solution.tex}\\\textcolor{DescText}{Teacher-facing worked interpretations and decisions for the ten C4 problems.}}, filenode, text width=16.0cm]
[{\code{G11\_T2\_C5\_catchUp.tex}\\\textcolor{DescText}{Student-facing C5 practice constructing mean intervals when long-run population spread is known.}}, filenode, text width=16.0cm]
[{\code{G11\_T2\_C5\_catchUp\_solution.tex}\\\textcolor{DescText}{Teacher-facing calculations and interpretations for the ten C5 interval problems.}}, filenode, text width=16.0cm]
[{\code{G11\_T2\_C6\_catchUp.tex}\\\textcolor{DescText}{Student-facing C6 practice planning sample sizes, interval widths, and precision upgrades.}}, filenode, text width=16.0cm]
[{\code{G11\_T2\_C6\_catchUp\_solution.tex}\\\textcolor{DescText}{Teacher-facing worked sample-size solutions for the twelve C6 planning scenarios.}}, filenode, text width=16.0cm]
[{\code{G11\_T2\_C7\_catchUp.tex}\\\textcolor{DescText}{Student-facing C7 practice building, comparing, and using confidence intervals for proportions.}}, filenode, text width=16.0cm]
[{\code{G11\_T2\_C7\_catchUp\_solution.tex}\\\textcolor{DescText}{Teacher-facing step-by-step solutions and recommendations for the C7 proportion problems.}}, filenode, text width=16.0cm]
]
\end{forest}
\end{center}
\end{landscape}
\begin{landscape}
\subsection{\code{4-extra\_material/}}
\begin{center}
\begin{forest}
for tree={grow=east, parent anchor=east, child anchor=west, anchor=west, edge={-Latex,draw=TreeLine}, l sep=6mm, s sep=1.2mm, font=\sffamily\fontsize{6.2}{7.2}\selectfont},
[{\textbf{4-extra\_material/}\\\textcolor{DescText}{Supplementary Lesson 1 assessment work and consolidated Lesson 3 worked material.}}, rootnode, text width=6.0cm
[{\code{G11\_T2\_L1\_supplementary.tex}\\\textcolor{DescText}{Student-facing extra Lesson 1 problems on known-spread mean intervals across confidence levels and samples.}}, filenode, text width=16.0cm]
[{\code{G11\_T2\_L1\_supplementary\_solution.tex}\\\textcolor{DescText}{Teacher-facing worked solutions to the supplementary Lesson 1 problem set.}}, filenode, text width=16.0cm]
[{\code{G11\_T2\_L3\_MATERIAL.tex}\\\textcolor{DescText}{Consolidated worked Lesson 3 material for C8 t-intervals, C9 t-versus-z comparison, and C10 mean conclusions.}}, filenode, text width=16.0cm]
]
\end{forest}
\end{center}
\end{landscape}

\section{Resource categories}

\begin{itemize}
\item \textbf{Student-facing pages:} the unsuffixed exams, activities, mock exams, project briefs, project problems, catch-up sheets, and supplementary sheet are printable LaTeX assessment or lesson sources.
\item \textbf{Teacher materials:} every \code{\_solution.tex} or \code{\_solutions.tex} source supplies worked answers. \code{G11\_T2\_L2\_reference\_midterm.tex}, the five practice sources, and \code{G11\_T2\_L3\_MATERIAL.tex} also contain worked solutions despite not consistently using a solution suffix.
\item \textbf{Lesson and assessment resources:} Lesson 1 introduces known-spread intervals; Lesson 2 develops criteria C6–C7 and midterm assessment; Lesson 3 develops unknown-variance t inference and conclusions (C8–C10). Catch-ups remediate C1–C7 individually.
\item \textbf{Stylesheets and JavaScript:} none are stored in this directory. The AI web-page project describes a deliverable but does not include its HTML, CSS, or JavaScript implementation.
\item \textbf{Images and media:} none are stored locally. Sixteen standalone assessment sheets reference the shared logo outside this directory.
\item \textbf{Data and configuration:} no separate datasets or configuration files are present. Numeric observations and grouped tables are embedded directly in the LaTeX sources; the Excel project only specifies CSV import as a student deliverable.
\item \textbf{Downloadable documents:} all 41 inventoried files are \code{.tex} source documents. There are no prebuilt PDF, Word, spreadsheet, or other downloadable binaries here.
\item \textbf{Generated/build-related files:} none are present or intentionally version-controlled in this tree.
\end{itemize}

\section{Important relationships}

\begin{itemize}
\item Twenty-five sources use a shared \code{preamble\_exams.tex} to provide document setup and custom commands such as \code{\textbackslash{}ExamTitleBlock}, \code{\textbackslash{}ExamSection}, and \code{ExamProblems}; the remaining sixteen define their own \code{article} preamble and directly reference the shared \code{logo.png}.
\item Both shared resources live outside \code{term\_2}, under \code{book2026-2027/preamble/}. The sixteen standalone sources consistently refer to the logo as \code{../../../preamble/logo.png} from their immediate subdirectories.
\item Nine shared-preamble sources use \code{../../../preamble/preamble\_exams.tex}, which resolves from their subdirectories to the repository's \code{book2026-2027/preamble/} directory. Sixteen instead use \code{../../preamble/preamble\_exams.tex}; from the source-file directory that points to a nonexistent \code{grade11/preamble/}. Whether those sixteen are compiled from a working directory that makes the reference valid \textbf{requires verification}.
\item Student/teacher pairs share the same problem contexts: Lesson 1 exam and practice, C6 and C7 activities, standard and hard mock midterms, project problems, each C1–C7 catch-up, and Lesson 1 supplementary work each connect to the correspondingly named solution source.
\item \code{G11\_T2\_L2\_reference\_midterm.tex} is a worked reference rather than the answer key for either named mock. \code{G11\_T2\_L2\_mock\_midterm\_solution.tex} and \code{G11\_T2\_L2\_mock\_midterm\_hard\_solution.tex} are the direct keys for those two mocks.
\item \code{G11\_T2\_L3\_MATERIAL.tex} consolidates worked C8, C9, and C10 content; the three Lesson 3 practice files split those criteria into separate documents. The C8/C9 practice is confidence-interval work, while C10 focuses on inference conclusions and decision errors.
\item The two Lesson 3 project briefs are alternative implementation tracks: one asks students to present an AI-generated web page, while the other asks for an Excel workbook with imported CSV observations, dashboard output, and macros. Neither implementation nor dataset is included here.
\item LaTeX compilation is the only build relationship visible in this directory. There are no HTML pages, stylesheets, scripts, package manifests, build scripts, or compiled outputs to connect locally.
\end{itemize}

\section{Inventory notes}

\begin{itemize}
\item \textbf{Documented:} 41 existing files and 4 directories (excluding the \code{term\_2} root and this inventory file).
\item \textbf{Excluded:} no transient or generated items were found; there were no \code{.DS\_Store}, editor-cache, LaTeX auxiliary, temporary, dependency, or compiled PDF files to omit.
\item \textbf{Requires verification:} the intended compilation working directory for the sixteen sources whose shared-preamble path is \code{../../preamble/preamble\_exams.tex}.
\end{itemize}
\end{document}
